\section{Additional Experiments (camera-ready)}
\label{app:b-additional}

This appendix collects analyses that supplement, but do not directly
test, the pre-registered stale-coherence hypothesis \textbf{H}. They
are reported here so the underlying telemetry is available to
readers while keeping the main-text experimental section focused on
the stale-coherence phenomenon.

\subsection{B.1 Decision-Token Efficiency}
\label{app:b1}

\paragraph{Hypothesis.} Spearman rank correlation
$\rho(\text{RAE-rank}, \text{DTE-rank}) < 0.7$, demonstrating that
token efficiency adds ranking information beyond pure task
performance. Specifically, 7B-tier models with frequent short
directives can outrank larger models with rare rich directives on
the DTE axis.

\paragraph{Design.} Reuses E1 logs (\S\ref{sec:exp:e1}) for the
cross-model comparison; adds 15 cadence-ablation runs (1 model
$\times$ 3 cadence levels $\times$ 5 seeds) to test whether DTE is
robust to the LLM's call-frequency knob. The DTE metric stack
(\texttt{dte\_per\_completion\_token},
\texttt{dte\_per\_decision}, \texttt{first\_token\_latency\_p50})
is defined in \S\ref{sec:metrics}. This analysis asks whether
completion-token cost is decoupled from task performance and is
reported only as supporting telemetry alongside the stale-coherence
phenomenon.

\paragraph{Expected result.} Spearman $\rho < 0.7$ between RAE rank
and DTE rank across the 3-model E1 matrix indicates that the DTE
axis provides additional ranking information beyond pure RAE.
